% ============================================================================
%  MAPA CONCEPTUAL — Introducción a la Lógica y Problemas Computacionales (ILyPC)
%  Pseudocódigo estilo cátedra (Alan Rodas Bonjour, UNAHUR)
%  Fuentes: NOTACION.md (reglas de cátedra) + Material de PC UNLP (Bordignon e Iglesias)
%  Compilar con: pdflatex mapa_ilypc.tex
% ============================================================================
\documentclass[a4paper,11pt,landscape]{article}

% --- Paquetes esenciales ---
\usepackage[utf8]{inputenc}
\usepackage[T1]{fontenc}
\usepackage[spanish]{babel}

% --- Matemática ---
\usepackage{amsmath,amssymb}

% --- Diseño ---
\usepackage[margin=1.4cm]{geometry}
\usepackage{array}
\usepackage{booktabs}

% --- Color y TikZ ---
\usepackage{xcolor}
\usepackage{tikz}
\usepackage{forest}
\usetikzlibrary{arrows.meta}

% --- Hipervínculos ---
\usepackage[hidelinks]{hyperref}

% ============================================================================
%  COLORES Y ESTILOS
% ============================================================================
\definecolor{azulT}{HTML}{1F4E79}      % teoría / algoritmo
\definecolor{verdeT}{HTML}{2E7D32}     % pseudocódigo
\definecolor{naranjaT}{HTML}{E65100}   % parámetros
\definecolor{rojoT}{HTML}{C62828}      % alternativas
\definecolor{violetaT}{HTML}{6A1B9A}   % repeticiones
\definecolor{grisT}{HTML}{546E7A}      % notas

\tikzset{
  nodoRaiz/.style={
    draw=azulT, very thick, fill=azulT!15,
    rounded corners=4pt, align=center, font=\bfseries\large\sffamily
  },
  nodoTema/.style={
    draw=azulT, thick, fill=azulT!10,
    rounded corners=4pt, align=center, font=\bfseries\small\sffamily
  },
  nodoAlgo/.style={
    draw=azulT, fill=azulT!5,
    rounded corners=3pt, align=center, font=\small\sffamily
  },
  nodoPseudo/.style={
    draw=verdeT, fill=verdeT!6,
    rounded corners=3pt, align=center, font=\small\sffamily
  },
  nodoParam/.style={
    draw=naranjaT, fill=naranjaT!6,
    rounded corners=3pt, align=center, font=\small\sffamily
  },
  nodoAlt/.style={
    draw=rojoT, fill=rojoT!6,
    rounded corners=3pt, align=center, font=\small\sffamily
  },
  nodoRep/.style={
    draw=violetaT, fill=violetaT!6,
    rounded corners=3pt, align=center, font=\small\sffamily
  },
  nodoNota/.style={
    draw=grisT, dashed, fill=grisT!6,
    rounded corners=3pt, align=center, font=\footnotesize\sffamily
  }
}

% ============================================================================
%  DOCUMENTO
% ============================================================================
\begin{document}

\begin{center}
  {\LARGE\bfseries\sffamily Mapa Conceptual}\\[2pt]
  {\large\sffamily Introducción a la Lógica y Problemas Computacionales (ILyPC)}\\[2pt]
  {\small\sffamily\color{grisT} Pseudocódigo estilo cátedra (UNAHUR) --- según \texttt{NOTACION.md} y material de PC (UNLP)}
\end{center}

\vspace{2mm}

% ---------------------------------------------------------------------------
%  MAPA CONCEPTUAL (árbol forest, raíz arriba, ramas en fila)
% ---------------------------------------------------------------------------
\begin{center}
\resizebox{\textwidth}{!}{
\begin{forest}
  for tree={
    draw, semithick,
    rounded corners=3pt,
    align=center,
    l sep=3.5mm,                % separación vertical entre niveles
    s sep=5mm,                  % separación horizontal entre ramas
    edge={-{Stealth[length=2mm]}, semithick, gray!60},
    font=\small\sffamily,
  }
  [ILyPC, nodoRaiz
    [Algoritmo\\{\small\itshape secuencia finita de pasos}, nodoTema
      [{Secuencial: se ejecuta\\una sola vez}, nodoAlgo]
      [{Con ciclos: instrucciones\\que se repiten}, nodoAlgo]
      [{PC: descomposición,\\patrones, abstracción}, nodoNota]
    ]
    [Pseudocódigo ILyPC\\{\small\itshape notación de la cátedra}, nodoTema
      [{Punto de entrada:\\\texttt{Al empezar a ejecutar}}, nodoPseudo]
      [{Procedimientos:\\\texttt{Definir} + infinitivo}, nodoPseudo]
      [{Pipes: \texttt{|} nivel 1,\\\texttt{||} nivel 2}, nodoPseudo]
      [{Primitivas: \texttt{Leer()},\\\texttt{Mostrar()}, \texttt{Mover}}, nodoPseudo]
      [{Sensores: \texttt{¿pregunta?}\\\mbox{$\to$} Sí/No}, nodoPseudo]
    ]
    [Usos de parámetros\\{\small\itshape en procedimientos}, nodoTema
      [{Entrada: dato que\\el procedimiento recibe}, nodoParam]
      [{Salida: dato que\\produce y devuelve}, nodoParam]
      [{Entrada y salida:\\recibe y modifica}, nodoParam]
      [{Reutilizar el módulo\\con datos distintos}, nodoNota]
    ]
    [Alternativas\\condicionales, nodoTema
      [{Simple:\\\texttt{Si ¿sensor? Entonces}\\{\small\ttfamily |acción}}, nodoAlt]
      [{Doble:\\\texttt{Si ... Entonces}\\{\small\ttfamily |acción}\\{\small\ttfamily Sino}\\{\small\ttfamily |otra acción}}, nodoAlt]
      [{No anidar \texttt{Si}:\\extraer procedimiento}, nodoNota]
    ]
    [Repeticiones\\{\small\itshape ciclos}, nodoTema
      [{Determinada:\\\texttt{Repetir N veces}\\{\small\ttfamily |acción}}, nodoRep]
      [{Indeterminada:\\\texttt{Mientras ¿condición?}\\{\small\ttfamily |acción}}, nodoRep]
      [{No anidar \texttt{Repetir}:\\extraer procedimiento}, nodoNota]
    ]
  ]
\end{forest}}
\end{center}

\vspace{3mm}

% ---------------------------------------------------------------------------
%  EJEMPLO CANÓNICO (del NOTACION.md de la cátedra)
% ---------------------------------------------------------------------------
\begin{center}
  \begin{minipage}{0.58\textwidth}
    \begin{center}
      {\small\bfseries\sffamily Ejemplo canónico: leer dos números y mostrar el mayor}
    \end{center}
    {\ttfamily\small
      Al empezar a ejecutar\\
      |Determinar el mayor de dos numeros ingresados
      \par\vspace{2pt}
      Definir Determinar el mayor de dos numeros ingresados\\
      |Leer(numeroA)\\
      |Leer(numeroB)\\
      |Comparar y mostrar mayor
      \par\vspace{2pt}
      Definir Comparar y mostrar mayor\\
      |Si ¿numeroA es mayor que numeroB? Entonces\\
      ||Mostrar(numeroA)\\
      |Sino\\
      ||Mostrar(numeroB)
    }
  \end{minipage}
  \hspace{4mm}
  \begin{minipage}{0.34\textwidth}
    \begin{center}
      {\small\bfseries\sffamily Leyenda de colores}
    \end{center}
    {\small\sffamily
    \begin{tabular}{@{}ll@{}}
      \color{azulT}{\rule{5mm}{2mm}} & Teoría y algoritmo\\
      \color{verdeT}{\rule{5mm}{2mm}} & Pseudocódigo ILyPC\\
      \color{naranjaT}{\rule{5mm}{2mm}} & Usos de parámetros\\
      \color{rojoT}{\rule{5mm}{2mm}} & Alternativas condicionales\\
      \color{violetaT}{\rule{5mm}{2mm}} & Repeticiones\\
      \color{grisT}{\rule{5mm}{2mm}} & Notas y reglas\\
    \end{tabular}
    }
    \vspace{2mm}

    {\footnotesize\sffamily\color{grisT}
      Reglas de oro: primitivas escritas tal cual las da el programa; verbos en infinitivo; no anidar \texttt{Si} ni \texttt{Repetir}.
    }
  \end{minipage}
\end{center}

\vspace{2mm}
\begin{center}
  {\footnotesize\sffamily\color{grisT}
  Fuentes: \texttt{NOTACION.md} (pseudocódigo estilo apuntes ILyPC, Alan Rodas Bonjour, UNAHUR) --- \textit{Introducción al Pensamiento Computacional} (Bordignon e Iglesias, UNLP/EDUCAR). El libro \textit{Introducción a la Lógica y Problemas Computacionales} no está en el material consultado.
  }
\end{center}

\end{document}
